\pdfbookmark[1]{Auszug aus dreijährigen täglichen Beobachtungen der magnetischen Deklination in Göttingen.}{dreijaehrige-declination}
\begin{center}\large\textit{Auszug aus dreijährigen täglichen Beobachtungen der magnetischen Deklination in Göttingen.}\end{center}
\begin{center}\small Resultate aus den Beobachtungen des magnetischen Vereins.\quad 1836.\quad III.\end{center}

Bei dem unaufhörlichen Wechsel kleinerer und größerer Schwankungen der magnetischen Deklination, die wir unregelmäßig nennen, insofern ihr Auftreten keiner zeitlichen Regel folgt, gibt es nur einen Weg, das Regelmäßige herauszulösen. Man muss sehr viele Beobachtungen nach einem bestimmten Plan anstellen, sie lange mit beharrlicher Konsequenz fortsetzen und sie geeignet kombinieren, um Mittelwerte zu gewinnen, in denen der Einfluss der stets auf die einzelnen Beobachtungen wirkenden Abweichungen so weit wie möglich verschwindet. Während der Vormittagsstunden nimmt die Deklination in unseren Gegenden gewöhnlich zu, aber an einem Tag stark und an einem anderen wenig. Zuweilen (wenn auch selten) beobachtet man sogar zu der Stunde, in der die Deklination gewöhnlich am größten ist, eine kleinere als in den frühen Stunden desselben Tages. Die Ursache der vormittäglichen Zunahme mag an jedem Tag wirksam sein. Ihre Wirkung wird aber durch andere, regellos dazwischentretende Kräfte zuweilen verstärkt, zuweilen abgeschwächt und zuweilen ganz verdeckt. Wie viel die regelmäßige Ursache tatsächlich bewirkt und wie unterschiedlich sie in den verschiedenen Jahreszeiten wirkt, lässt sich nicht an einzelnen oder wenigen Tagen, sondern nur an Mittelwerten aus sehr vielen Tagen
\SourceLeaf{563}erkennen. Ähnlich verhält es sich mit den allmählichen Änderungen, die zumindest sehr lange in derselben Richtung fortschreiten. Wir nennen sie säkular, weil eine lange Reihe von Jahren nötig ist, damit sie sich zu vielen Grad summieren. Einzelne Beobachtungen, die nur wenige Jahre auseinanderliegen, können uns darüber noch keine sichere Auskunft geben, selbst wenn sie am gleichen Monatstag und zur gleichen Stunde angestellt werden. Konsequentes Beobachten liefert dagegen Mittelwerte, mit denen wir schon nach wenigen Jahren vorwegnehmen können, was sich sonst mit einiger Annäherung erst nach mehreren Jahrzehnten feststellen ließe.

Von diesem Gesichtspunkt aus habe ich von Anfang an die tägliche Bestimmung der absoluten Deklination, immer zu denselben Stunden, in die Beobachtungen am hiesigen magnetischen Observatorium aufgenommen. Solche Arbeiten erhalten ihren Wert erst dadurch, dass man sie lange und ohne Unterbrechung fortsetzt. Um eher darauf rechnen zu können, dass dies möglich ist, wollte ich lieber zunächst einen beschränkten Plan wählen, als auf einmal zu viel erfassen. Deshalb werden täglich nur zwei Bestimmungen vorgenommen, vormittags um 8 Uhr und nachmittags um 1 Uhr mittlerer Zeit. Diese Stunden ließen sich am leichtesten mit anderen Pflichten vereinbaren. Für sie sprach auch, dass bei regelmäßigem Verlauf der magnetischen Bewegungen die Nadel um 1 Uhr nachmittags immer nahe am Maximum der Deklination steht und um 8 Uhr vormittags während des größeren Teils des Jahres nahe am Minimum. Beobachtungen zu bestimmten Stunden \textit{wahrer} Sonnenzeit wären an sich noch etwas natürlicher gewesen. Hier, wo es hauptsächlich auf die konsequente Durchführung nach demselben Prinzip ankam, musste jedoch die viel größere Bequemlichkeit einer Einteilung nach mittlerer Zeit den Ausschlag geben.

Diese regelmäßigen Aufzeichnungen begannen am ersten Januar 1834. Die ersten zweieinhalb Monate sind jedoch vom folgenden Auszug ausgeschlossen, weil der Aufhängungsfaden in dieser Zeit öfter aufgewickelt werden musste und sich dadurch der Nullpunkt der Torsion änderte. Darauf achtete man anfangs nicht genug. Vom 17. März an wurde ein stärkerer Aufhängungsfaden (aus zweihundert Einzelfäden) verwendet, nachdem sein Torsionsnullpunkt zuvor genau berichtigt worden war. Wann immer man später an diesem Faden oder an einem anderen Umstand, der mit den
\SourceLeaf{564}Reduktionselementen zusammenhing, etwas änderte, nahm man die nötigen Berichtigungen oder Änderungen der Reduktionselemente vor. In den ersten Monaten teilten sich verschiedene hinreichend geübte Beobachter mit mir die Beobachtungen. Seit dem 1. Oktober 1834 stellt sie jedoch regelmäßig Herr Dr.~\textsc{Goldschmidt} an. Nur wenn er verhindert ist, vertreten ihn andere geübte Beobachter.

Die monatlichen Mittelwerte dieser Bestimmungen bis Januar 1835 habe ich bereits in den \textit{Göttingischen gelehrten Anzeigen} von 1834 und 1835 [S.\ 519 und S.\ 528 dieses Bandes] mitgeteilt. Hier folgen nun die Mittelwerte für drei vollständige Jahrgänge.

\Needspace{31\baselineskip}
\begin{center}\textit{Mittelwert der westlichen magnetischen Deklination in Göttingen.}\end{center}
\begin{center}\small
\begin{tabular}{rl|rrr|rrr}
&&\multicolumn{3}{c|}{8 Uhr vorm.}&\multicolumn{3}{c}{1 Uhr nachm.}\\\hline
1834&März, zweite Hälfte&$18^0$&$38'$&$16''0$&$18^0$&$46'$&$40''4$\\
&April&&36&6.9&&47&3.8\\
&Mai&&36&28.2&&47&15.4\\
&Juni&&37&40.7&&47&59.5\\
&Juli&&37&57.5&&48&19.0\\
&August&&38&48.1&&49&11.0\\
&September&&36&58.4&&46&32.3\\
&Oktober&&37&18.4&&44&47.2\\
&November&&37&38.4&&43&4.3\\
&Dezember&&37&54.8&&41&32.7\\[4pt]
1835&Januar&&37&51.5&&42&14.4\\
&Februar&&37&3.5&&42&29.4\\
&März&&34&47.5&&44&55.2\\
&April&&32&57.7&&46&31.6\\
&Mai&&32&13.4&&45&17.1\\
&Juni&&32&56.4&&44&41.3\\
&Juli&&34&8.0&&44&42.8\\
&August&&34&12.4&&46&56.8\\
&September&&33&21.2&&44&27.6\\
&Oktober&&33&23.0&&43&5.3\\
&November&&36&15.3&&43&49.5\\
&Dezember&&35&25.9&&40&19.1\\[4pt]
1836&Januar&&35&2.4&&40&34.6\\
&Februar&&33&26.7&&41&15.2\\
&März&&31&1.4&&43&16.4\\
&April&&26&32.9&&43&42.6\\
&Mai&&28&0.8&&44&37.2
\end{tabular}
\end{center}
\SourceLeaf{565}\Needspace{15\baselineskip}
\begin{center}\small
\begin{tabular}{rl|rrr|rrr}
&&\multicolumn{3}{c|}{8 Uhr vorm.}&\multicolumn{3}{c}{1 Uhr nachm.}\\\hline
1836&Juni&$18^0$&$27'$&$35''1$&$18^0$&$42'$&$52''4$\\
&Juli&&26&54.2&&42&26.0\\
&August&&25&42.4&&41&45.0\\
&September&&26&14.6&&40&59.6\\
&Oktober&&27&34.0&&40&32.8\\
&November&&29&21.0&&36&54.3\\
&Dezember&&29&13.7&&35&46.8\\[4pt]
1837&Januar&&27&35.3&&37&46.2\\
&Februar&&27&35.6&&36&28.3\\
&März&&25&44.2&&39&4.2
\end{tabular}
\end{center}

Hier mögen nun einige Kombinationen dieser Beobachtungen folgen.

Der Unterschied zwischen der Deklination am Vormittag und am Nachmittag hat in den Mittelwerten durchweg dasselbe Vorzeichen. Wie die Größe dieses Unterschieds von der Jahreszeit abhängt, zeigt die folgende Übersicht.
\begin{center}\small
\begin{tabular}{l||rr|rr|rr||rr}
&\multicolumn{2}{c|}{1834. 1835}&\multicolumn{2}{c|}{1835. 1836}&\multicolumn{2}{c||}{1836. 1837}&\multicolumn{2}{c}{Mittel}\\\hline\hline
April&$10'$&$56''9$&$13'$&$33''9$&$17'$&$9''7$&$13'$&$53''5$\\
Mai&10&47.2&13&3.7&16&36.4&13&29.1\\
Juni&10&18.8&11&44.9&15&17.3&12&27.0\\
Juli&10&21.5&10&34.8&15&31.8&12&9.4\\
August&10&22.9&12&44.4&16&2.6&13&3.3\\
September&9&33.9&11&6.4&14&45.0&11&48.4\\
Oktober&7&28.8&9&42.3&12&58.8&10&3.3\\
November&5&25.9&7&34.2&7&33.3&6&51.1\\
Dezember&3&37.9&4&53.2&6&33.1&5&1.4\\
Januar&4&22.9&5&32.2&10&10.9&6&42.0\\
Februar&5&25.9&7&48.5&8&52.7&7&22.4\\
März&10&7.7&12&15.0&13&20.0&11&54.2\\\hline\hline
Mittel&$8'$&$14''2$&$10'$&$2''8$&$12'$&$54''3$&$10'$&$23''8$
\end{tabular}
\end{center}
Man sieht, dass der Unterschied sowohl in den Mittelwerten als auch in jedem einzelnen Jahr im Dezember am kleinsten war. Das erscheint auch ganz natürlich, da die mit der Tageszeit wechselnden Änderungen notwendig einer Einwirkung der Sonne zugeschrieben werden müssen, auch wenn wir noch nicht wissen, \textit{wie} diese Einwirkung zustande kommt. Dagegen mag es zunächst auffallen, dass die ungleich größeren Unterschiede der Sommermonate nicht um die Zeit der
\SourceLeaf{566}Sonnenwende am größten waren, sondern im Juni und Juli kleiner als im April, Mai und August. Das gilt besonders, weil die Übereinstimmung aller drei einzelnen Jahre in diesem Punkt vermuten lässt, dass es kein Zufall ist. Dabei darf man jedoch nicht übersehen, dass in den der Sonnenwende nächsten Monaten das Minimum der Deklination schon auf eine frühere Stunde fällt und die gesamte Zunahme daher merklich größer wäre als die von 8 Uhr an gerechnete Bewegung.

Es fällt ferner auf, dass der Unterschied im zweiten Jahr in jedem einzelnen Monat größer war als im ersten und im dritten wiederum größer als im zweiten. Die Unterschiede sind aber viel zu groß, als dass man darin etwas sehen dürfte, was auf eine säkulare Zunahme hinausläuft. Vielmehr ist zu erwarten, dass bei einer Fortsetzung der Beobachtungen über mehrere Jahre Schwankungen in beiden Richtungen nicht ausbleiben. Auf jeden Fall lernen wir daraus, dass sich auch in der Einwirkung der Sonne auf den Erdmagnetismus ein Jahr von einem anderen abheben kann, etwa so, wie sich ein ganzer Sommer oder ein ganzer Winter durch seine Witterung erheblich von anderen unterscheidet. Eben deshalb wird man die Mittelwerte erst durch mehrjährige Beobachtungen genau bestimmen können.

Dass an einzelnen Tagen der Unterschied zwischen Vormittags- und Nachmittagsdeklination ausnahmsweise das entgegengesetzte Vorzeichen haben kann, wurde schon oben bemerkt. Wie selten solche Ausnahmen sind, zeigt sich daran, dass während der dreijährigen Beobachtungen nur vierzehn Fälle dieser Art vorkamen, im Durchschnitt also einer unter 79 Tagen. Ich führe sie hier an und gebe jeweils an, um wie viel die Deklination um 8 Uhr morgens größer war als um 1 Uhr nachmittags.
\begin{center}\small
\begin{tabular}{rlr|rr||rlr|rr}
1834&August&15&$6'$&$8''0$&1835&November&8&$3'$&$42''2$\\
&Dezember&24&3&43.0&&Dezember&8&18&35.6\\
&Dezember&25&0&38.2&1836&Januar&20&0&46.3\\
&Dezember&26&2&20.3&&Juli&20&5&8.8\\
1835&Januar&30&0&23.8&&November&9&11&9.5\\
&Februar&7&0&32.5&1837&Februar&13&4&1.0\\
&Oktober&4&0&43.1&&März&14&1&22.6
\end{tabular}
\end{center}
Dass von diesen vierzehn Ausnahmen zwölf in die Wintermonate und nur zwei in die Sommermonate fallen, ist ganz natürlich, da die geringe
\SourceLeaf{567}regelmäßige Sonnenwirkung in den Wintermonaten leichter von einer anomalen Bewegung übertroffen werden kann als die viel größere in den Sommermonaten.

Um zu prüfen, inwieweit die vorliegenden Beobachtungen die säkulare Änderung schon erkennen lassen, wurden die Monatsmittel des ersten Jahres mit den entsprechenden Mitteln des zweiten und ebenso die des zweiten mit denen des dritten verglichen. Von den 48 Vergleichen, die sich auf diese Weise ergeben (der unvollständige März 1834 ist hiervon ebenso wie von den übrigen Kombinationen ausgeschlossen), ergeben 47 eine Abnahme und nur einer eine Zunahme. Diese ist deshalb in der folgenden Übersicht mit dem Minuszeichen bezeichnet.
\Needspace{17\baselineskip}
\begin{center}\small\setlength{\tabcolsep}{3pt}
\textit{Jährliche Abnahme der Deklination.}\par\smallskip
\begin{tabular}{l||rr|rr||rr|rr||rr}
&\multicolumn{4}{c||}{Erstes Jahr}&\multicolumn{4}{c||}{Zweites Jahr}&\multicolumn{2}{c}{Mittel}\\
&\multicolumn{2}{c|}{8 Uhr vorm.}&\multicolumn{2}{c||}{1 Uhr nachm.}&\multicolumn{2}{c|}{8 Uhr vorm.}&\multicolumn{2}{c||}{1 Uhr nachm.}&&\\\hline\hline
April&$3'$&$9''2$&$0'$&$32''2$&$6'$&$24''8$&$2'$&$49''0$&$3'$&$13''8$\\
Mai&4&14.8&1&58.3&4&12.6&0&39.9&2&46.4\\
Juni&4&44.3&3&18.2&5&21.3&1&48.9&3&48.1\\
Juli&3&49.5&3&36.2&7&13.8&2&16.8&4&14.1\\
August&4&35.7&2&14.2&8&30.0&5&11.8&5&7.9\\
September&3&37.2&2&4.7&7&6.6&3&28.0&4&4.1\\
Oktober&3&55.4&1&41.9&5&49.0&2&32.5&3&29.6\\
November&1&23.1&$-0$&45.2&6&54.3&6&55.2&3&36.8\\
Dezember&2&28.9&1&13.6&6&12.2&4&32.3&3&36.7\\
Januar&2&49.1&1&39.8&7&27.1&2&48.4&3&41.1\\
Februar&3&36.8&1&14.2&5&51.1&4&46.9&3&52.2\\
März&3&46.1&1&38.8&5&17.2&4&12.2&3&43.6\\\hline\hline
Mittel&3&30.8&1&42.2&6&21.7&3&20.2&3&46.2
\end{tabular}
\end{center}
Dass der Vergleich der Vormittagsmittel hier meist eine stärkere Abnahme ergibt als der Vergleich der Nachmittagsmittel, ist nur eine andere Formulierung der obigen Beobachtung, dass die täglichen Änderungen im ersten Jahr geringer waren als im zweiten und im zweiten geringer als im dritten. Jener Unterschied ist daher lediglich als zufällig anzusehen, und bei einer längeren Fortsetzung der Beobachtungen ist auch ein Unterschied im entgegengesetzten Sinn zu erwarten. Sofern es also keinen hinreichenden Grund gibt, dem einen Ergebnis den Vorzug vor dem anderen zu geben, bleibt nur, sich an das Mittel aus
\SourceLeaf{568}beiden zu halten. Dieses Mittel beträgt für das erste Jahr $2'36''5$ und für das zweite $4'55''9$. Man könnte versucht sein, darin einen Beweis dafür zu sehen, dass sich die Abnahme der Deklination beschleunigt. Das wäre jedoch nur eine schlechte Begründung für eine an sich richtige Sache. Bekanntlich hat die Deklination, die im vorigen Jahrhundert in ganz Europa zunahm, im gegenwärtigen Jahrhundert ihr Maximum erreicht und geht seitdem wieder zurück. Der Natur der Sache nach muss dieser Übergang eine zunächst unmerkliche und nach und nach stärker werdende Abnahme erzeugen. Mangels früherer Beobachtungen lässt sich das Jahr, in dem dieser Übergang für Göttingen stattfand, zwar nicht mit Bestimmtheit angeben. Nach den von anderen Orten bekannt gewordenen Beobachtungen muss man es aber für beträchtlich weiter zurückliegend halten, als sich aus jenen beiden Zahlen ergäbe, wenn man sie als reine Wirkungen der langsamen Bewegung ansehen wollte, die wir Säkularbewegung nennen. Ebenso wenig lässt sich nach allen sonstigen Erfahrungen eine so starke Änderung wie $2'19''4$ als regelmäßige Zunahme für ein Jahr gelten lassen. Wir halten daher auch diesen Unterschied größtenteils für zufällig, sodass vorerst, bis weiter reichende Erfahrungen vorliegen, das Mittel $3'46''2$ als einjährige Abnahme der Deklination für 1834---1837 gelten muss.

Da der Unterschied zwischen der Deklination zur Vormittags- und zur Nachmittagsstunde einer so offenkundig mit der Jahreszeit wechselnden Ungleichheit unterliegt, stellt sich die Frage, ob nur eine der beiden, vorwiegend eine oder beide zugleich an einer jahreszeitlich bedingten Änderung beteiligt sind und welche Gesetze ihr zugrunde liegen. Um diese Gesetze zu ermitteln, wird eine längere Reihe von Jahren noch notwendiger sein als für den bloßen Unterschied der Deklinationen. Inzwischen wird man aber gern sehen, was die bisherigen Beobachtungen, soweit sie reichen, aussagen.

Zu diesem Zweck wurden zunächst die Mittelwerte aus je zwölf Monaten für die drei Beobachtungsjahre berechnet. Sie lauten wie folgt.
\begin{center}\small
\begin{tabular}{l|rrr|rrr}
&\multicolumn{3}{c|}{8 Uhr vorm.}&\multicolumn{3}{c}{1 Uhr nachm.}\\\hline
1834---1835&$18^0$&$37'$&$12''5$&$18^0$&$45'$&$27''0$\\
1835---1836&&33&42.0&&43&44.8\\
1836---1837&&27&20.3&&40&14.6
\end{tabular}
\end{center}

\SourceLeaf{569}Diese Mittelwerte sind als gültig für den mittleren Tag jedes Rechnungsjahres anzusehen, die ersten also für den 1. Oktober 1834 und so weiter.

Der Vergleich der einzelnen Monate jedes Jahres mit dem zugehörigen Mittelwert ergibt die folgenden Unterschiede.
\begin{center}\small
\textit{Deklination um 8 Uhr vormittags.}\par\smallskip
\begin{tabular}{l|rr|rr|rr||rr}
&\multicolumn{2}{c|}{Erstes Jahr}&\multicolumn{2}{c|}{Zweites Jahr}&\multicolumn{2}{c||}{Drittes Jahr}&\multicolumn{2}{c}{Mittel}\\\hline
April&$-1'$&$5''9$&$-0'$&$44''3$&$-0'$&$47''4$&$-0'$&$52''5$\\
Mai&$-0$&44.6&$-1$&28.6&$+0$&40.5&$-0$&30.9\\
Juni&$+0$&27.9&$-0$&45.6&$+0$&14.8&$-0$&1.0\\
Juli&$+0$&44.7&$+0$&26.0&$-0$&26.1&$+0$&14.9\\
August&$+1$&35.3&$+0$&30.4&$-1$&37.9&$+0$&9.3\\
September&$-0$&14.4&$-0$&20.8&$-1$&5.7&$-0$&33.6\\
Oktober&$+0$&5.6&$-0$&19.0&$+0$&13.7&$-0$&0.1\\
November&$+0$&25.6&$+2$&33.3&$+2$&0.7&$+1$&39.9\\
Dezember&$+0$&42.0&$+1$&43.9&$+1$&53.4&$+1$&26.4\\
Januar&$+0$&38.7&$+1$&20.4&$+0$&15.0&$+0$&44.7\\
Februar&$-0$&9.3&$-0$&15.3&$+0$&15.3&$-0$&3.1\\
März&$-2$&25.3&$-2$&40.6&$-1$&36.1&$-2$&14.0
\end{tabular}
\end{center}
\Needspace{16\baselineskip}
\begin{center}\small
\textit{Deklination um 1 Uhr nachmittags.}\par\smallskip
\begin{tabular}{l|rr|rr|rr||rr}
&\multicolumn{2}{c|}{Erstes Jahr}&\multicolumn{2}{c|}{Zweites Jahr}&\multicolumn{2}{c||}{Drittes Jahr}&\multicolumn{2}{c}{Mittel}\\\hline
April&$+1'$&$36''8$&$+2'$&$46''8$&$+3'$&$28''0$&$+2'$&$37''2$\\
Mai&$+1$&48.4&$+1$&32.3&$+4$&22.6&$+2$&34.4\\
Juni&$+2$&32.5&$+0$&56.5&$+2$&37.8&$+2$&2.3\\
Juli&$+2$&52.0&$+0$&58.0&$+2$&11.4&$+2$&0.5\\
August&$+3$&44.0&$+3$&12.0&$+1$&30.4&$+2$&48.8\\
September&$+1$&5.3&$+0$&42.8&$+0$&45.0&$+0$&51.0\\
Oktober&$-0$&39.8&$-0$&39.5&$+0$&18.2&$-0$&20.4\\
November&$-2$&22.7&$+0$&4.7&$-3$&20.3&$-1$&52.8\\
Dezember&$-3$&54.3&$-3$&25.7&$-4$&27.8&$-3$&55.9\\
Januar&$-3$&12.6&$-3$&10.2&$-2$&28.4&$-2$&57.1\\
Februar&$-2$&57.6&$-2$&29.6&$-3$&46.3&$-3$&4.5\\
März&$-0$&31.8&$-0$&28.4&$-1$&10.4&$-0$&43.5
\end{tabular}
\end{center}
Als Mittelwerte aus drei Jahren sind die Zahlen der letzten Spalte einigermaßen, wenn auch erst sehr unvollkommen, vom Einfluss der unregelmäßigen Anomalien befreit. Sie enthalten aber offenbar noch die säkulare Änderung. Um sie herauszurechnen, muss ihr Betrag zwischen der
\SourceLeaf{570}Mitte jedes Monats und dem 1. Oktober für die ersten sechs Monate mit negativem und für die letzten sechs mit positivem Vorzeichen angebracht werden. Legen wir den oben bestimmten Zwölfmonatswert $3'46''2$ zugrunde, so erhalten wir die folgenden Ergebnisse.
\begin{center}\small
\begin{tabular}{l|rr|rr||rr}
&\multicolumn{2}{c|}{8 Uhr vorm.}&\multicolumn{2}{c||}{1 Uhr nachm.}&\multicolumn{2}{c}{Mittel}\\\hline
April&$-2'$&$35''6$&$+0'$&$54''2$&$-0'$&$50''7$\\
Mai&$-1$&55.3&$+1$&10.0&$-0$&22.6\\
Juni&$-1$&6\phantom{.}6&$+0$&56.7&$-0$&4.9\\
Juli&$-0$&32.0&$+1$&13.6&$+0$&20.8\\
August&$-0$&18.8&$+2$&20.7&$+1$&0.9\\
September&$-0$&43.0&$+0$&41.6&$-0$&0.7\\
Oktober&$+0$&9.3&$-0$&11.0&$-0$&0.8\\
November&$+2$&8.0&$-1$&24.7&$+0$&21.6\\
Dezember&$+2$&13.3&$-3$&9.0&$-0$&27.8\\
Januar&$+1$&50.3&$-1$&51.5&$-0$&0.6\\
Februar&$+1$&21.3&$-1$&40.1&$-0$&9.4\\
März&$-0$&30.9&$+0$&59.6&$+0$&14.3
\end{tabular}
\end{center}
In diesen Ergebnissen zeigt sich schon so viel Regelmäßigkeit, wie man von nur dreijährigen Beobachtungen erwarten konnte. Die erste Spalte zeigt, um wie viel die Vormittagsdeklination in den einzelnen Monaten von der mittleren Vormittagsdeklination abweicht. Ebenso gibt die zweite Spalte den Unterschied zwischen der Nachmittagsdeklination jedes Monats und der mittleren Nachmittagsdeklination an. Dabei muss man bedenken, dass die letztere selbst um $10'23''8$ größer ist als die mittlere Vormittagsdeklination.

Bemerkenswert scheint nun, dass in allen zwölf Monaten die Vormittags- und die Nachmittagsdeklination \textit{auf entgegengesetzten Seiten} über ihre Mittelwerte hinausschwanken. In den fünf Wintermonaten von Oktober bis Februar ist die Vormittagsdeklination größer als ihr Mittelwert und die Nachmittagsdeklination kleiner. Beide Umstände tragen also \textit{zugleich} dazu bei, in dieser Jahreszeit den gesamten Unterschied unter seinen Mittelwert zu bringen. In den übrigen sieben Monaten tritt gerade das Gegenteil ein. Überdies sind diese entgegengesetzten Schwankungen durchschnittlich beinahe gleich groß. Daher heben sie sich in ihrem Mittelwert, den die letzte Spalte darstellt, fast auf. Anders ausgedrückt enthält das Mittel zwischen der magnetischen Deklination um 8 Uhr vormittags und um 1 Uhr nachmittags
\SourceLeaf{571}neben den unregelmäßigen Anomalien und der säkularen Abnahme keine erheblichen von der Jahreszeit abhängigen Schwankungen. Zumindest tritt kein Unterschied zwischen Sommer- und Wintermonaten mit Sicherheit hervor.

Der Mittelwert selbst, aus sämtlichen dreijährigen Beobachtungen abgeleitet, wäre für den 1. Oktober 1835
\[
=18^0\;37'\;56''9
\]
anzusetzen. Es versteht sich übrigens von selbst, dass hier nur der Mittelwert aus den für unsere Beobachtungen gewählten Stunden gemeint ist. Der Mittelwert aus \textit{allen} Stunden des Tages könnte davon etwas abweichen, wenn auch wahrscheinlich nur wenig. Alle bisherigen Untersuchungen zeigen aber zur Genüge, dass sich darüber ohne sehr langwierige Arbeiten nichts mit Sicherheit feststellen lässt.

Bisher war nur von den monatlichen Mittelwerten die Rede. Den vollständigen Abdruck der einzelnen Beobachtungen hielt man vorerst für überflüssig. Solange sie nur von einem Ort vorliegen, können sie nur insofern von Interesse sein, als sich das unregelmäßige Hin- und Herspringen an ihnen erkennen lässt. Dieses Ziel erreicht man jedoch besser als durch das bloße Betrachten der Zahlen, wenn man sie methodisch kombiniert. So führt man die Größe der Schwankungen auf ein bestimmtes Maß zurück und kann den allgemeinen Charakter verschiedener Zeiträume hinsichtlich ihrer größeren oder kleineren Schwankungen genau vergleichen. Der Kürze halber verstehe ich hier unter der Schwankung der magnetischen Deklination die Differenz zu ihrem Wert am Vortag zur selben Stunde und (in Analogie zu den sogenannten mittleren Beobachtungsfehlern) unter der mittleren Schwankung während eines beliebigen Zeitraums die Quadratwurzel aus dem Mittel der Quadrate der einzelnen Schwankungen. Dabei ist zu beachten, dass man, wenn mehrere gleiche oder als gleich betrachtete Zeiträume nachträglich zu einem einzigen vereinigt werden sollen, zur Bestimmung des Gesamtmittels nicht das arithmetische Mittel aus den einzelnen mittleren Schwankungen nehmen darf. Man muss zuerst auf deren Quadrate zurückgehen, daraus das arithmetische Mittel bilden und sich an dessen Quadratwurzel halten. Die Ergebnisse der so über die dreijährigen
\SourceLeaf{572}Beobachtungen angestellten Rechnung enthält die folgende Tabelle in Bogensekunden.
\Needspace{19\baselineskip}
\begin{center}\small
\textit{Mittlere Schwankung der magnetischen Deklination während der drei Jahre}\par
1834---1837.\par\smallskip
\begin{tabular}{l|r|r|r||r||r|r|r||r}
&\multicolumn{4}{c||}{8 Uhr vorm.}&\multicolumn{4}{c}{1 Uhr nachm.}\\\cline{2-9}
&I&II&III&Mittel&I&II&III&Mittel\\\hline
April&74&126&205&147&129&101&264&180\\
Mai&192&124&277&207&158&183&210&185\\
Juni&172&171&199&181&95&151&217&162\\
Juli&213&243&287&250&119&184&252&193\\
August&264&253&269&262&175&165&307&225\\
September&162&325&207&241&172&143&161&159\\
Oktober&116&296&216&222&182&202&242&210\\
November&79&205&308&218&170&173&126&158\\
Dezember&132&324&71&206&184&206&154&182\\
Januar&146&274&138&196&174&212&154&181\\
Februar&116&146&164&143&178&183&129&165\\
März&100&109&366&228&127&153&246&183\\\hline
Mittel&157&229&238&211&156&174&213&183
\end{tabular}
\end{center}

Von den einzelnen Beobachtungen seien hier noch die \textit{größten} Schwankungen angeführt, die im Laufe der drei Jahre bei der Vormittags- und der Nachmittagsdeklination vorkamen. Die Vormittagsdeklination war am 8. Oktober 1835 um $20'1''$ größer als am 7. Oktober und die Nachmittagsdeklination am 24. April 1836 um $13'0''$ größer als am vorhergehenden Tag. Umgekehrt ist auch öfter vorgekommen, dass die Vormittags- oder die Nachmittagsdeklination an zwei aufeinanderfolgenden Tagen genau gleich war. Bei den monatlichen mittleren Schwankungen rücken diese Extreme natürlich viel näher zusammen. Die großen Unterschiede zwischen den einzelnen Monaten bleiben in dieser Hinsicht dennoch sehr bemerkenswert. Nach der obigen Übersicht betrug die mittlere Schwankung der Vormittagsdeklination im März 1837 $6'6''$, im Dezember 1836 dagegen nur $1'11''$.

Ob im Allgemeinen zu einer Tageszeit größere Schwankungen vorherrschen als zu einer anderen, lässt sich aus den Ergebnissen für unsere beiden Stunden noch nicht mit Sicherheit entscheiden. Im ersten Jahr sind die beiden Mittelwerte
\SourceLeaf{573}nahezu gleich. In den beiden anderen Jahren überwiegen die Vormittagsschwankungen. Der Unterschied zwischen den Endergebnissen für alle drei Jahre, $3'31''$ und $3'3''$, ist aber zu klein, um ihn aufgrund so weniger Jahre für gesichert zu halten, obwohl in den Mittelwerten für die einzelnen Monate in der vierten und achten Spalte zehn Monate einen Unterschied in derselben Richtung ergeben haben.

Fasst man Vormittags- und Nachmittagsbeobachtungen zusammen, so erhält man die folgenden mittleren Schwankungen.
\begin{center}\small
\begin{tabular}{l|r|r|r||r}
&Jahr I&Jahr II&Jahr III&Mittel\\\hline
April&108&114&237&164\\
Mai&176&156&245&196\\
Juni&139&161&208&172\\
Juli&173&215&270&223\\
August&224&214&289&244\\
September&167&251&185&204\\
Oktober&152&254&229&216\\
November&133&190&235&191\\
Dezember&160&271&120&195\\
Januar&160&245&146&189\\
Februar&150&166&148&155\\
März&114&133&312&206
\end{tabular}
\end{center}
\Needspace{7\baselineskip}
\begin{center}\small
\textit{Mittelwerte.}\par\smallskip
\begin{tabular}{l|r|r|r||r}
&Jahr I&Jahr II&Jahr III&Mittel\\\hline
Juli---Dezember&170&234&228&213\\
Übrige Monate&143&167&223&181\\
Ganzes Jahr&158&204&226&198
\end{tabular}
\end{center}

Nach den Zahlen der vierten Spalte sind die Schwankungen in den Monaten Juli bis Dezember etwas größer als in den übrigen sechs Monaten. Die Mittelwerte $3'33''$ und $3'1''$ sind aber wohl zu wenig verschieden, um mit Sicherheit darauf schließen zu können, dass diese Jahreszeit größere Schwankungen stärker begünstigt. Das gilt besonders, weil der Unterschied vor allem in dem einen Jahr 1835---1836 so stark hervorgetreten ist.

Sehr deutlich ist dagegen die unterschiedliche Veränderlichkeit im Vergleich der drei einzelnen Jahre. Der Mittelwert für das dritte Jahr ist fast um die Hälfte größer als der Mittelwert für das erste. Das Gesamtmittel
\SourceLeaf{574}aus sämtlichen bisherigen Beobachtungen, $3'18''$, könnte sich daher bei längerer Fortsetzung noch erheblich ändern.

Dies sind die Ergebnisse, die sich aus den bisherigen täglichen Aufzeichnungen der magnetischen Deklination ziehen lassen. Es ist sehr zu wünschen, dass ähnliche Arbeiten an mehreren Orten ausgeführt werden. An einigen hat man seit Kurzem bereits damit begonnen. Wenn die Beobachtungen wie in Mailand nicht nach der Ortszeit, sondern genau gleichzeitig mit den hiesigen vorgenommen werden, bietet der Vergleich der einzelnen Tage noch Gelegenheit zu anderen Kombinationen, die. sobald sie eine etwas beträchtliche Zeit umfassen können, von großem Interesse sein werden. Beobachter, die es ähnlich halten, also ihre Aufzeichnungen zu Zeiten machen, die mit den hiesigen übereinstimmen, werden daher gebeten, die Ergebnisse jedes einzelnen Tages mitzuteilen. Es reicht dabei aber aus, sie nur in Skalenteilen anzugeben. Die Umrechnung in Bogenteile kann dann entfallen, sofern man zugleich die nötigen Reduktionselemente angibt.

